\chapter{Experimental: reduction and stratification}
\label{chap:experimental-reduction}

\texttt{Lean4Lean/Experimental/} holds sixteen research modules that are not in the dependency cone
of anything shipping. Four of them (the logical-relation line) are Chapter~\ref{chap:experimental-logrel};
the remaining twelve, 6\,321 lines, are this chapter. They fall into two clusters. The
\emph{\texttt{VExpr} cluster} works over the real model syntax and is an earlier, freestanding copy of
the Church-Rosser and stratification material: \srcloc{Lean4Lean/Experimental/NormalEq.lean}{1} (513
lines) and \srcloc{Lean4Lean/Experimental/ParallelReduction.lean}{1} (913 lines), both headed
\texttt{TODO: remove, this is now part of ChurchRosser.lean}, plus
\srcloc{Lean4Lean/Experimental/Stratified.lean}{1} (340 lines),
\srcloc{Lean4Lean/Experimental/StratifiedUntyped.lean}{1} (324 lines) and
\srcloc{Lean4Lean/Experimental/Stronger.lean}{1} (628 lines). The \emph{\texttt{SExpr} and
normalization cluster} works over a separate quotiented syntax and over five unfinished semantic
models: \srcloc{Lean4Lean/Experimental/SExpr.lean}{1} (1\,324 lines),
\srcloc{Lean4Lean/Experimental/UniqueTyping.lean}{1} (266),
\srcloc{Lean4Lean/Experimental/StepIndexed.lean}{1} (88),
\srcloc{Lean4Lean/Experimental/MoreStepIndexed.lean}{1} (489),
\srcloc{Lean4Lean/Experimental/DomainTheory.lean}{1} (242),
\srcloc{Lean4Lean/Experimental/Thierry.lean}{1} (363) and
\srcloc{Lean4Lean/Experimental/Thierry2.lean}{1} (831).

\emph{Thesis correspondence.} The \texttt{VExpr} cluster is a direct formalization of
\texttt{unique.tex}: \texttt{NormalEq} is the proof-irrelevance relation $\equiv_p$ of
\S\texttt{sec:kappa}, \texttt{ParRed} and \texttt{CParRed} are $\gg_\kappa$ and $\ggg_\kappa$ of
\S\texttt{sec:church\_rosser}, and \texttt{HasType1}/\texttt{IsDefEq1} are one layer of the
stratification $\vdash_n$ of \S\texttt{sec:unique} that the thesis uses to break the circularity
between unique typing and Church-Rosser. The shipping counterpart of the first two files is
\texttt{Theory/Typing/ChurchRosser.lean} (Chapter~\ref{chap:metatheory}); the shipping route to
unique typing is different (\texttt{HasTypeStratified}, not $\vdash_n$), so
\texttt{Stratified.lean} has no shipping counterpart at all. The normalization cluster corresponds
to \texttt{normalization.tex}, which is a 22-line stub ending in the word \texttt{UNFINISHED}; there
is accordingly nothing to check these files against.

\emph{What the branches contributed.} \texttt{trproj} changed nothing here. \texttt{iota} changed 28
lines across five files: one field plus one supporting import in \texttt{NormalEq.lean}, eight
lines in \texttt{ParallelReduction.lean}, nine in \texttt{Stratified.lean}, eight in
\texttt{StratifiedUntyped.lean} and one in \texttt{Stronger.lean}. Every one of them is the minimum
needed to keep an unrelated experiment compiling after \texttt{VEnv} grew a \texttt{pats} field
(\ref{struct:venv}) and \texttt{VEnv.IsDefEq} grew a \texttt{pat} constructor
(\ref{rule:isdefeq-pat}). \texttt{Lean4Lean.Experimental} is not a \texttt{lake} default target, but
the CI workflow builds it explicitly, which is what forces the maintenance.

\emph{Assessment.} The one place in this group where real design judgement was required is the
stratification of the $\iota$ rule's premises in \texttt{IsDefEq1} (\ref{inductive:expb-isdefeq1}),
and it was made correctly: the typing premise goes to the abstract \texttt{HasType1} parameter, as
\texttt{beta}/\texttt{eta}/\texttt{proofIrrel} do, while the side-condition premises stay recursive,
as \texttt{appDF}/\texttt{trans} do; the pair \ref{thm:expb-induction1} and
\ref{thm:expb-isdefeq1-induction} shows the new rule round-trips. Everything else falls into three
groups: the \texttt{Typing.pat\_env} field and the \texttt{pat} case of
\texttt{VEnv.IsDefEqU.church\_rosser} (\ref{thm:expb-isdefequ-church-rosser}) are verbatim duplicates
of code the branch already wrote in \texttt{ChurchRosser.lean}; \texttt{StratifiedUntyped.lean}
repeats the same stratification decision a second time for the type-erased judgment rather than
making a fresh one; and \texttt{VEnv.IsDefEq.toTyping}'s \texttt{pat} case
(\ref{thm:expb-todyping}) and the \texttt{pats \_ \_ := False} stub in \texttt{Stronger.lean}, a file
that dead-ends, are neither a duplicate nor a repeated decision, just the minimum glue needed to keep
each file compiling. None of the 28 lines carries a comment, and
almost all of them land on ground that was already unproved before the branch: this chapter contains
one file-local \texttt{sorry} in each of \texttt{Stratified.lean} and \texttt{StratifiedUntyped.lean},
two in \texttt{ParallelReduction.lean}, four in \texttt{Stronger.lean} and thirty in
\texttt{SExpr.lean}, all of them master. Two files, \texttt{Thierry.lean} and \texttt{Thierry2.lean},
declare \texttt{axiom mySorry : $\alpha$} and are therefore inconsistent by construction.

\section{The abstract typing interface and proof irrelevance}

\begin{definition}[Pointwise definitionally-equal contexts, \texttt{IsDefEqCtx}]
\label{inductive:expb-isdefeqctx}
\lean{Lean4Lean.IsDefEqCtx}
\leanok
\srcloc{Lean4Lean/Experimental/NormalEq.lean}{11}
For a fixed relation \texttt{IsDefEqU} and prefix $\Gamma_0$, \texttt{IsDefEqCtx IsDefEqU} $\Gamma_0\,
\Gamma_1\,\Gamma_2$ is generated by \texttt{zero}, which relates $\Gamma_0$ to itself, and
\texttt{succ}, which extends a related pair $\Gamma_1,\Gamma_2$ by hypotheses $A_1, A_2$ with
\texttt{IsDefEqU} $\Gamma_1\,A_1\,A_2$. It is the context-conversion relation used by the
\texttt{DFC} lemmas below.
\end{definition}

\begin{definition}[Abstract typing interface \texttt{Typing}, with the $\iota$-rule field]
\label{structure:expb-typing}
\lean{Lean4Lean.Typing, Lean4Lean.Typing.pat_env, Lean4Lean.Typing.HasType}
\leanok
\uses{ind:vexpr, struct:venv, inductive:pattern-core, inductive:pattern-rhs, inductive:expb-isdefeqctx}
\contrib{mixed}\srcloc{Lean4Lean/Experimental/NormalEq.lean}{20}
\thesisref{unique.tex \S\texttt{sec:kappa} and \S\texttt{sec:church\_rosser}: the properties assumed
of $\vdash_n$ before the Church-Rosser argument is run.}
\texttt{Typing} is a record bundling an environment \texttt{env}, a universe count, a typed relation
$\Gamma\vdash e_1\equiv e_2 : A$ and its erased form $\Gamma\vdash e_1\equiv e_2$, together with
every property the Church-Rosser argument needs: congruence and structural rules, $\beta$, $\eta$,
proof irrelevance, weakening and its inverse, substitution, level instantiation, context conversion
(\texttt{isDefEq\_DFC}), regularity (\texttt{has\_type}, \texttt{is\_type}), head inversions
(\texttt{bvar\_inv}, \texttt{sort\_inv}, \texttt{const\_inv}, \texttt{forallE\_inv},
\texttt{app\_inv}, \texttt{lam\_inv}), unique typing (\texttt{uniq}) and definitional inversion
(\texttt{univ\_defInv}, \texttt{forallE\_defInv}). It also carries an abstract family
\texttt{Pat p r} of reduction rules with their regularity (\texttt{pat\_wf}), uniqueness
(\texttt{pat\_uniq}), neutrality (\texttt{pat\_onArgs}) and a realization
\texttt{extra\_pat} of definitional axioms by patterns. The \texttt{iota} branch adds the field
\texttt{pat\_env : env.pats p r} $\to$ \texttt{Pat p r}, saying that every reduction rule registered
in the environment is one of the abstract rules. \texttt{Typing.HasType} abbreviates the reflexive
instance of the typed relation.
\reviewnote{No term of type \texttt{Lean4Lean.Typing} exists anywhere in the repository, so the new
obligation is never discharged here; its real counterpart is the class field
\texttt{VEnv.Params.pat\_env} (\ref{axiom:cr-params-pat-env}), which is discharged with
\texttt{pat\_env := id}. Unlike that one, this field has no docstring.}
\end{definition}

\begin{lemma}[Pattern-side helpers for the interface]
\label{family:expb-normaleq-pattern-aux}
\lean{Lean4Lean.Typing.symm_ctx, Lean4Lean.Typing.IsDefEqU.weakN, Lean4Lean.Typing.IsDefEq.weakN, Lean4Lean.Typing.IsDefEqU.instN, Lean4Lean.Typing.IsDefEq.instN, Lean4Lean.Pattern.Check.OK.weakN, Lean4Lean.Pattern.Check.OK.instN, Lean4Lean.Typing.IsDefEqU.apply_pat, Lean4Lean.Pattern.Matches.hasType}
\leanok
\uses{structure:expb-typing, inductive:pattern-rhs, inductive:pattern-matches}
\srcloc{Lean4Lean/Experimental/NormalEq.lean}{104}
\texttt{symm\_ctx} symmetrises \texttt{IsDefEqCtx}; the four \texttt{weakN}/\texttt{instN} lemmas
unpack the corresponding interface fields; \texttt{Check.OK.weakN} and \texttt{Check.OK.instN}
transport the side-condition predicate of a pattern rule through weakening and substitution;
\texttt{IsDefEqU.apply\_pat} says that instantiating a rule right-hand side at pointwise
definitionally equal metavariable assignments gives definitionally equal results; and
\texttt{Matches.hasType} says every metavariable captured by a match of a well-typed term is itself
typeable.
\end{lemma}
\begin{proof}
\uses{structure:expb-typing}
\leanok
The transport lemmas are \texttt{Check.OK.map} applied to the weakening/substitution fields, after
rewriting with \texttt{Pattern.RHS.liftN\_apply} and \texttt{Pattern.RHS.instN\_apply};
\texttt{apply\_pat} and \texttt{Matches.hasType} are inductions on the right-hand side resp.\ the
match, both using the interface's \texttt{app\_inv}.
\end{proof}

\begin{definition}[The proof-irrelevance relation, \texttt{NormalEq}]
\label{inductive:expb-normaleq}
\lean{Lean4Lean.NormalEq}
\leanok
\uses{structure:expb-typing}
\srcloc{Lean4Lean/Experimental/NormalEq.lean}{165}
\thesisref{unique.tex \S\texttt{sec:kappa}, the rule block for $\Gamma\vdash e\equiv_p e'$.}
\texttt{NormalEq TY} $\Gamma\,e_1\,e_2$, written $\Gamma\vdash e_1\equiv_p e_2$, is generated by
reflexivity at a well-typed term, level congruence at sorts and at constants, congruence for
application (with both sides typed at a common $\forall A.\,B$), congruence for $\lambda$ and
$\forall$ where the two binder types are only required to be definitionally equal to a common $A$,
the two one-sided $\eta$ rules \texttt{etaL} and \texttt{etaR}, and proof irrelevance. It has no
symmetry, transitivity or reduction rule: those are theorems.
\end{definition}

\begin{lemma}[Proof irrelevance implies definitional equality]
\label{thm:expb-normaleq-defeq}
\lean{Lean4Lean.NormalEq.defeq}
\leanok
\uses{inductive:expb-normaleq, structure:expb-typing}
\srcloc{Lean4Lean/Experimental/NormalEq.lean}{202}
\thesisref{unique.tex, the ``regularity of reductions'' lemma, item 3.}
If $\Gamma\vdash e_1\equiv_p e_2$ then $\Gamma\vdash e_1\equiv e_2$ in the erased relation of the
interface.
\end{lemma}
\begin{proof}
\uses{structure:expb-typing}
\leanok
Induction on the derivation, replaying each rule with the matching interface field. The
\texttt{lamDF} and \texttt{forallEDF} cases retype both binders at the common $A$ using
\texttt{TY.trans} and \texttt{TY.symm}; the two $\eta$ cases combine \texttt{TY.lamDF} with
\texttt{TY.eta} after recovering the domain from \texttt{TY.is\_type} and
\texttt{TY.forallE\_inv}.
\end{proof}

\begin{lemma}[Structural metatheory of proof irrelevance]
\label{family:expb-normaleq-struct}
\lean{Lean4Lean.NormalEq.weakN, Lean4Lean.NormalEq.instN, Lean4Lean.NormalEq.instN_r, Lean4Lean.NormalEq.defeqDFC, Lean4Lean.NormalEq.defeq_l, Lean4Lean.NormalEq.weakN_inv_DFC, Lean4Lean.NormalEq.weakN_iff, Lean4Lean.NormalEq.apply_pat}
\leanok
\uses{inductive:expb-normaleq, structure:expb-typing, inductive:expb-isdefeqctx}
\srcloc{Lean4Lean/Experimental/NormalEq.lean}{239}
\thesisref{unique.tex, the ``regularity of reductions'' lemma, substitution item.}
$\equiv_p$ is stable under weakening (\texttt{weakN}) and reflects it
(\texttt{weakN\_inv\_DFC}, whence the equivalence \texttt{weakN\_iff}), is stable under substitution
on either side (\texttt{instN}, \texttt{instN\_r}), transports along a definitionally equal context
(\texttt{defeqDFC}, \texttt{defeq\_l}), and is a congruence for rule right-hand sides
(\texttt{apply\_pat}).
\end{lemma}
\begin{proof}
\uses{structure:expb-typing, thm:expb-normaleq-defeq}
\leanok
Each is an induction over the derivation. The two $\eta$ cases need
\texttt{VExpr.lift\_liftN'} to re-associate the lifts. Weakening inversion is the hard one: it
generalises simultaneously over a lifting and a context conversion and uses \texttt{TY.uniq} and
\texttt{TY.forallE\_defInv} to recover the common types that lifting destroys.
\end{proof}

\begin{theorem}[Proof irrelevance is an equivalence relation]
\label{thm:expb-normaleq-trans}
\lean{Lean4Lean.NormalEq.trans, Lean4Lean.NormalEq.symm}
\leanok
\uses{inductive:expb-normaleq}
\srcloc{Lean4Lean/Experimental/NormalEq.lean}{451}
\thesisref{unique.tex, the ``regularity of reductions'' lemma, item 2.}
$\equiv_p$ is symmetric and transitive; with the \texttt{refl} rule it is an equivalence relation on
well-typed terms.
\end{theorem}
\begin{proof}
\uses{family:expb-normaleq-struct, thm:expb-normaleq-defeq, structure:expb-typing}
\leanok
Symmetry is a direct induction that swaps the two sides, using \texttt{TY.trans} in the
\texttt{forallEDF} case to move the binder equations across. Transitivity is a well-founded
recursion on \texttt{meas} $e_1 +$ \texttt{meas} $e_2 +$ \texttt{meas} $e_3$, where the private
measure weights a $\lambda$ by $3$ so that the $\eta$-versus-$\eta$ clash decreases; that case
inverts the inner transitive result and closes with \texttt{NormalEq.weakN\_iff}, and the mixed
$\lambda$-versus-$\eta$ cases retype the abstraction with \texttt{TY.forallEDF}.
\end{proof}

\section{Parallel reduction and Church-Rosser}

\begin{definition}[Parallel reduction, \texttt{ParRed}]
\label{inductive:expb-parred}
\lean{Lean4Lean.ParRed}
\leanok
\uses{structure:expb-typing, inductive:pattern-matches, inductive:pattern-rhs}
\srcloc{Lean4Lean/Experimental/ParallelReduction.lean}{18}
\thesisref{unique.tex \S\texttt{sec:church\_rosser}, the rule block for $\gg_\kappa$; the $\iota$
and $K^+$ rules of the thesis are folded into the single \texttt{extra} rule.}
$\Gamma\vdash e\gg_\kappa e'$ is reflexive on atoms and congruent in every subterm, with a
$\beta$ rule reducing $(\lambda x{:}A.\,e_1)\,e_2$ to $e_1'[e_2'/x]$ after reducing both parts, and
an \texttt{extra} rule: given an abstract rule \texttt{TY.Pat p r}, a match \texttt{p.Matches}
$e\,m_1\,m_2$ whose side conditions \texttt{r.2.OK} hold in the interface, and a parallel reduction
$m_2\,a \gg_\kappa m_2'\,a$ of every captured metavariable, $e$ reduces to
\texttt{r.1.apply} $m_1\,m_2'$.
\end{definition}

\begin{definition}[Reducible terms, \texttt{NonNeutral}]
\label{def:expb-nonneutral}
\lean{Lean4Lean.NonNeutral}
\leanok
\uses{structure:expb-typing}
\srcloc{Lean4Lean/Experimental/ParallelReduction.lean}{29}
\thesisref{unique.tex \S\texttt{sec:church\_rosser}: ``the compatibility rules only apply if none of
the substantive rules are applicable''.}
\texttt{NonNeutral TY} $\Gamma\,e$ holds when $e$ is syntactically a $\beta$ redex, or when some
abstract rule matches $e$ with its side conditions satisfied.
\end{definition}

\begin{definition}[Complete parallel reduction, \texttt{CParRed}]
\label{inductive:expb-cparred}
\lean{Lean4Lean.CParRed}
\leanok
\uses{inductive:expb-parred, def:expb-nonneutral}
\srcloc{Lean4Lean/Experimental/ParallelReduction.lean}{33}
\thesisref{unique.tex \S\texttt{sec:church\_rosser}, the definition of $\ggg_\kappa$.}
$\ggg_\kappa$ has the same rules as $\gg_\kappa$ except that the congruence rules for constants and
for applications carry the guard $\lnot$\texttt{NonNeutral}, so that a substantive rule is always
preferred. Note that the $\lambda$ and $\forall$ congruences are unguarded, and that two overlapping
pattern rules can still both fire, so the relation is deterministic only up to the interface's
\texttt{pat\_uniq}.
\end{definition}

\begin{lemma}[Basic properties of the two reductions]
\label{family:expb-parred-struct}
\lean{Lean4Lean.ParRed.rfl, Lean4Lean.ParRed.weakN, Lean4Lean.ParRed.instN, Lean4Lean.ParRed.defeq, Lean4Lean.ParRed.hasType, Lean4Lean.ParRed.defeqDFC, Lean4Lean.ParRed.apply_pat, Lean4Lean.ParRed.weakN_inv, Lean4Lean.CParRed.toParRed, Lean4Lean.CParRed.exists, Lean4Lean.Pattern.RHS.apply_liftN}
\leanok
\uses{inductive:expb-parred, inductive:expb-cparred, structure:expb-typing}
\srcloc{Lean4Lean/Experimental/ParallelReduction.lean}{48}
\thesisref{unique.tex, the ``properties of $\gg$'' lemma (the reflexivity, defeq, substitution,
$\ggg\!\Rightarrow\!\gg$ and existence items).}
$\gg_\kappa$ is reflexive on all terms, stable under weakening, substitution and context conversion,
reflects weakening, is contained in the interface's definitional equality, and preserves typing.
Complete reduction refines parallel reduction, and every well-typed term has a complete reduct.
\end{lemma}
\begin{proof}
\uses{structure:expb-typing, def:expb-nonneutral}
\leanok
Each statement is an induction on the reduction, the \texttt{extra} cases going through
\texttt{Pattern.RHS.apply\_liftN} and the interface's \texttt{pat\_wf}.
\texttt{CParRed.exists} is instead a structural recursion on the term with a classical case split on
\texttt{NonNeutral}.
\end{proof}

\begin{theorem}[Triangle lemma]
\label{thm:expb-parred-triangle}
\lean{Lean4Lean.ParRed.triangle}
\leanok
\uses{inductive:expb-parred, inductive:expb-cparred, inductive:expb-normaleq}
\srcloc{Lean4Lean/Experimental/ParallelReduction.lean}{282}
\thesisref{unique.tex, the triangle lemma.}
If $\Gamma\vdash e : A$, $\Gamma\vdash e\gg_\kappa e'$ and $\Gamma\vdash e\ggg_\kappa e^\bullet$,
then there is $e^\circ$ with $\Gamma\vdash e'\gg_\kappa e^\circ$ and
$\Gamma\vdash e^\circ\equiv_p e^\bullet$.
\end{theorem}
\begin{proof}
\uses{family:expb-parred-struct, family:expb-normaleq-struct, structure:expb-typing}
\leanok
Well-founded recursion on the subject term, then induction on the complete reduction with inversion
on the parallel reduction. The $\beta$ cases use \texttt{ParRed.instN}; the cases where both sides
fire a pattern rule are aligned by the interface's \texttt{pat\_uniq}, and the cases where only one
fires are excluded by the \texttt{NonNeutral} guards. Roughly 150 lines, the hardest proof in the
file.
\end{proof}

\begin{theorem}[Compatibility of parallel reduction with proof irrelevance]
\label{thm:expb-normaleq-parred}
\lean{Lean4Lean.NormalEq.parRed, Lean4Lean.NormalEq.parRedS}
\uses{inductive:expb-normaleq, inductive:expb-parred}
\leanok
\srcloc{Lean4Lean/Experimental/ParallelReduction.lean}{687}
\thesisref{unique.tex, the compatibility lemma for $\gg_\kappa$ and $\equiv_p$.}
If $\Gamma\vdash e_1\equiv_p e_3$ and $\Gamma\vdash e_3\gg_\kappa e_2$, then there is $e_4$ with
$\Gamma\vdash e_1\gg_\kappa^{*} e_4$ and $\Gamma\vdash e_4\equiv_p e_2$. \texttt{parRedS} iterates
this along the reflexive-transitive closure.
\end{theorem}
\begin{proof}
\uses{family:expb-normaleq-struct, family:expb-parred-struct}
\textbf{Not proved in Lean.} Two \texttt{sorry}s remain, at lines 699 and 718: exactly the cases
where the reduction fires an \texttt{extra} (pattern) rule against a \texttt{constDF} resp.\ an
\texttt{appDF} step of $\equiv_p$. These are the cases the thesis discharges by observing that the
eliminator is small and the subject is therefore a proof. The proved part is an induction on the
$\equiv_p$ derivation with inversion on the reduction; the $\beta$-versus-$\eta$ clash goes through
\texttt{ParRedExt.parRed\_beta}, an auxiliary about reducing under a stack of applications and
lifts (\ref{thm:parredext-parred-beta} is the shipping copy).
\axfoot{direct sorry: ParallelReduction.lean:699 and :718}
\end{proof}

\begin{theorem}[Church-Rosser up to proof irrelevance]
\label{thm:expb-parred-church-rosser}
\lean{Lean4Lean.ParRed.church_rosser, Lean4Lean.ParRedS.church_rosser, Lean4Lean.Typing.CRDefEq}
\leanok
\uses{inductive:expb-parred, inductive:expb-normaleq, structure:expb-typing}
\srcloc{Lean4Lean/Experimental/ParallelReduction.lean}{431}
\thesisref{unique.tex, the Church-Rosser theorem.}
If $e$ is well-typed and parallel-reduces in one step to $e_1$ and to $e_2$, then $e_1$ and $e_2$
have further one-step reducts related by $\equiv_p$. \texttt{Typing.CRDefEq} $\Gamma\,e_1\,e_2$
packages the iterated form: both sides are typeable and reduce, along $\gg_\kappa^{*}$, to
$\equiv_p$-related terms; \texttt{ParRedS.church\_rosser} (line 799) establishes it from a common
ancestor.
\end{theorem}
\begin{proof}
\uses{thm:expb-parred-triangle, thm:expb-normaleq-parred, family:expb-parred-struct}
\leanok
The one-step case is the triangle lemma applied twice at the complete reduct supplied by
\texttt{CParRed.exists}, gluing the two results with transitivity and symmetry of $\equiv_p$. The
iterated case is a double induction along the two closures, commuting $\equiv_p$ past a further
reduction with \texttt{NormalEq.parRed} at each step, which is what makes the whole theorem
conditional.
\axfoot{sorryAx via NormalEq.parRed (ParallelReduction.lean:699, :718)}
\end{proof}

\begin{lemma}[\texttt{CRDefEq} is an equivalence contained in definitional equality]
\label{family:expb-crdefeq}
\lean{Lean4Lean.Typing.CRDefEq.normalEq, Lean4Lean.Typing.CRDefEq.refl, Lean4Lean.Typing.CRDefEq.defeq, Lean4Lean.Typing.CRDefEq.symm, Lean4Lean.Typing.CRDefEq.trans}
\leanok
\uses{thm:expb-parred-church-rosser, structure:expb-typing}
\srcloc{Lean4Lean/Experimental/ParallelReduction.lean}{819}
\texttt{CRDefEq} contains $\equiv_p$ and reflexivity at a well-typed term, is symmetric and
transitive, and implies the interface's definitional equality.
\end{lemma}
\begin{proof}
\uses{thm:expb-normaleq-parred, thm:expb-parred-church-rosser}
\leanok
Reflexivity, symmetry and \texttt{defeq} are unpacking (the last one chains
\texttt{ParRedS.defeq} on both sides with \texttt{NormalEq.defeq} in the middle). Transitivity glues
the two witnesses by applying \texttt{ParRedS.church\_rosser} to the shared middle term and then
commuting both $\equiv_p$ witnesses past the resulting reductions with \texttt{NormalEq.parRedS}.
\axfoot{sorryAx via NormalEq.parRed (ParallelReduction.lean:699, :718)}
\end{proof}

\begin{lemma}[Every concrete derivation is an interface derivation]
\label{thm:expb-todyping}
\lean{Lean4Lean.VEnv.IsDefEq.toTyping}
\leanok
\uses{structure:expb-typing, def:isdefeq}
\contrib{mixed}\srcloc{Lean4Lean/Experimental/ParallelReduction.lean}{839}
\thesisref{no thesis counterpart; this is the bridge from the concrete judgment to the abstract
interface.}
If $\Gamma\vdash e_1\equiv e_2 : A$ holds in \texttt{TY.env.IsDefEq}, then $e_1$ and $e_2$ are
related by \texttt{TY.IsDefEqU} and $e_1$ has type $A$ in the interface.
\end{lemma}
\begin{proof}
\uses{rule:isdefeq-pat, thm:pattern-realizes-took, structure:expb-typing}
\leanok
Induction on the derivation, replaying each rule through the corresponding field. The
\texttt{iota} branch supplies the new \texttt{pat} case (lines 857--859): the rule's realizer
premise is converted to the \texttt{Check.OK} form the interface expects by
\texttt{Pattern.Check.Realizes.toOK} applied to the induction hypotheses of the side conditions, the
environment rule is promoted to an abstract one by the new \texttt{pat\_env} field, and
\texttt{pat\_wf} concludes. The second component reuses the redex's typing, which is correct because
\texttt{IsDefEq.pat} types the reduct at the redex's type.
\reviewnote{This bridging step has no counterpart in \texttt{Theory/Typing/ChurchRosser.lean}: that
file proves its results directly about \texttt{env.IsDefEq} under a \texttt{[Params]} instance and
never needs to convert into a separate abstract record. The genuine duplicate of \texttt{iota}'s
pattern-matching plumbing is the \texttt{pat} case of \texttt{VEnv.IsDefEqU.church\_rosser} below
(\ref{thm:expb-isdefequ-church-rosser}), which mirrors \texttt{IsDefEq.church\_rosser} in
\texttt{ChurchRosser.lean} (lines 1390--1394) line for line.}
\end{proof}

\begin{theorem}[Church-Rosser for the concrete definitional equality]
\label{thm:expb-isdefequ-church-rosser}
\lean{Lean4Lean.VEnv.IsDefEqU.church_rosser}
\leanok
\uses{def:isdefeq, structure:expb-typing}
\contrib{mixed}\srcloc{Lean4Lean/Experimental/ParallelReduction.lean}{861}
\thesisref{unique.tex, the Church-Rosser theorem, transported to the concrete judgment.}
Every derivation of $\Gamma\vdash e_1\equiv e_2 : A$ in \texttt{VEnv.IsDefEq} yields
\texttt{TY.CRDefEq} $\Gamma\,e_1\,e_2$, that is, $e_1$ and $e_2$ reduce along $\gg_\kappa^{*}$ to
$\equiv_p$-related terms. The shipping counterpart is \ref{thm:isdefeq-church-rosser}.
\end{theorem}
\begin{proof}
\uses{thm:expb-todyping, family:expb-crdefeq, thm:expb-parred-church-rosser, rule:isdefeq-pat, thm:pattern-realizes-took}
\leanok
Induction on the derivation: the congruence rules use the \texttt{CRDefEq} closure lemmas, while
$\beta$, $\eta$ and \texttt{extra} exhibit an explicit one-step reduction. The \texttt{iota} branch
adds the \texttt{pat} case (lines 909--913), the only place in the file where an $\iota$ step has to
be realised as an actual reduction: it takes the witness to be \texttt{ParRed.extra} firing the very
rule \texttt{hp} (promoted by \texttt{pat\_env}, with side conditions converted by
\texttt{Realizes.toOK}) with every metavariable reduced by \texttt{.rfl}, and closes with
reflexivity of $\equiv_p$ at the reduct. That is exactly the shape of the pre-existing
\texttt{extra} case three lines above.
\axfoot{sorryAx via NormalEq.parRed (ParallelReduction.lean:699, :718)}
\end{proof}

\section{Stratified judgments}

\begin{definition}[Definitional inversion predicate, \texttt{DefInv}]
\label{def:expb-definv}
\lean{Lean4Lean.VEnv.DefInv, Lean4Lean.VEnv.DefInv.symm}
\leanok
\uses{def:isdefeq}
\srcloc{Lean4Lean/Experimental/Stratified.lean}{9}
\thesisref{unique.tex, the definition of definitional inversion preceding the unique-typing
theorem.}
\texttt{DefInv env U} $\Gamma\,e_1\,e_2$ is a head-shape compatibility condition on two types: two
$\forall$'s must have definitionally equal domains and, in the extended context, codomains; two
sorts must have equivalent levels; a sort against a $\forall$ is \texttt{False}; every other pair is
unconstrained. \texttt{DefInv.symm} shows the predicate is symmetric over an ordered environment,
moving the codomain equation across the domain equation with \texttt{IsDefEq.defeqDF\_l}. Nothing in
this file consumes \texttt{DefInv}: its intended consumers are in the commented-out block at lines
266--293.
\end{definition}

\begin{definition}[One typing layer, \texttt{HasType1}]
\label{inductive:expb-hastype1}
\lean{Lean4Lean.VEnv.HasType1}
\leanok
\uses{struct:venv}
\srcloc{Lean4Lean/Experimental/Stratified.lean}{30}
\thesisref{unique.tex \S\texttt{sec:unique}, the definition of $\Gamma\vdash_n e : \alpha$.}
\texttt{HasType1 env uvars IsDefEq1} is the typing judgment whose conversion rule appeals to an
\emph{abstract} parameter \texttt{IsDefEq1} rather than to the real definitional equality: variable,
constant, sort, application, abstraction, product, and a \texttt{defeq} rule taking
\texttt{IsDefEq1} $\Gamma\,A\,B\,(\mathtt{.sort}\ u)$. Each such parameter marks one alternation
between typing and conversion.
\end{definition}

\begin{definition}[One definitional-equality layer, \texttt{IsDefEq1}, with the $\iota$ rule]
\label{inductive:expb-isdefeq1}
\lean{Lean4Lean.VEnv.IsDefEq1, Lean4Lean.VEnv.IsDefEq1.pat}
\leanok
\uses{inductive:expb-hastype1, struct:venv, inductive:pattern-matches, def:pattern-check-realizes}
\contrib{mixed}\srcloc{Lean4Lean/Experimental/Stratified.lean}{46}
\thesisref{unique.tex \S\texttt{sec:unique} (the stratification) together with the $\iota$ rule of
\S\texttt{sec:kappa}.}
\texttt{IsDefEq1 env uvars HasType1 defEq} is definitional equality in which the typing premises
appeal to the abstract \texttt{HasType1} and the conversion rule \texttt{defeqDF} appeals to the
abstract \texttt{defEq}, while congruence, symmetry and transitivity recurse. The \texttt{iota}
branch adds the constructor \texttt{pat} (lines 71--75): from \texttt{env.pats p r}, a match
\texttt{p.Matches} $e\,m_1\,m_2$, an abstract typing $\Gamma\vdash e : A$, a realizer
\texttt{r.2.Realizes} $m_1\,m_2\,\mathit{chk}$ and a recursive $\Gamma\vdash t.1\equiv t.2.1 :
t.2.2$ for every triple $t$ in \textit{chk}, conclude $\Gamma\vdash e\equiv
\mathtt{r.1.apply}\,m_1\,m_2 : A$.
\reviewnote{The stratification choices are the delicate part and they are right: the typing premise
is an alternation point and goes to \texttt{HasType1}, as in \texttt{beta}/\texttt{eta}/
\texttt{proofIrrel}, while the side conditions are not and stay recursive, as in
\texttt{appDF}/\texttt{trans}. The \texttt{Realizes}-plus-list encoding rather than
\texttt{Check.OK} is forced by strict positivity. Neither point is commented, and the core rule this
mirrors (\ref{rule:isdefeq-pat}) does carry a docstring.}
\end{definition}

\begin{theorem}[Every derivation decomposes into one layer]
\label{thm:expb-induction1}
\lean{Lean4Lean.VEnv.IsDefEq.induction1}
\leanok
\uses{inductive:expb-isdefeq1, inductive:expb-hastype1, struct:orderedstrong}
\contrib{mixed}\srcloc{Lean4Lean/Experimental/Stratified.lean}{80}
\thesisref{unique.tex \S\texttt{sec:unique}, the ``$n$-provability basics'' lemma, items 3 and 4.}
Let \texttt{env} be \texttt{OrderedStrong}, $\Gamma$ well-formed, and let \texttt{hasType} and
\texttt{defEq} be relations closed under the one-layer judgments (hypotheses \texttt{hty} and
\texttt{hdf}). Then any $\Gamma\vdash e_1\equiv e_2 : A$ yields \texttt{HasType1} derivations for
$e_1$ and $e_2$ at $A$ together with an \texttt{IsDefEq1} derivation relating them.
\end{theorem}
\begin{proof}
\uses{thm:isdefeq-strong, inductive:expb-isdefeq1}
\textbf{Not proved in Lean.} One \texttt{sorry} remains at line 96, in the \texttt{constDF} case,
where a level-instantiation congruence for \texttt{ci.type} is needed; it predates the branch.
Otherwise the proof strengthens the hypothesis with \texttt{IsDefEq.strong}
(\ref{thm:isdefeq-strong}) and inducts on the strong derivation, rebuilding one layer per rule. The
\texttt{iota} branch adds the \texttt{pat} case (lines 121--122): the strong \texttt{pat} rule
already carries typings of both redex and reduct, so the two \texttt{HasType1} components are the
first projections of the two induction hypotheses, and the third is \texttt{IsDefEq1.pat} with
\texttt{hty} applied to the redex typing and the side-condition induction hypotheses.
\axfoot{direct sorry: Stratified.lean:96 (constDF)}
\reviewnote{Line 79 silently strengthens the exported hypothesis from \texttt{Ordered env} to
\texttt{OrderedStrong env}. The change is forced, because on this branch \texttt{IsDefEq.strong}
needs subject reduction for the registered $\iota$ rules (\ref{def:patsstrongon}), and
\ref{thm:wf-orderedstrong} shows no realistic instance is lost; but no comment says so, and a reader
diffing against master sees a weaker theorem with no rationale.}
\end{proof}

\begin{lemma}[One layer collapses back into the real judgment]
\label{thm:expb-isdefeq1-induction}
\lean{Lean4Lean.VEnv.IsDefEq1.induction, Lean4Lean.VEnv.HasType1.induction}
\leanok
\uses{inductive:expb-isdefeq1, inductive:expb-hastype1, def:isdefeq}
\contrib{mixed}\srcloc{Lean4Lean/Experimental/Stratified.lean}{142}
\thesisref{unique.tex \S\texttt{sec:unique}, the converse direction of the same lemma.}
If the abstract parameters are sound for the real judgments (\texttt{hty}, \texttt{hdf}, resp.\
\texttt{IH}), then \texttt{HasType1} implies \texttt{env.HasType} and \texttt{IsDefEq1} implies
\texttt{env.IsDefEq}. Together with \ref{thm:expb-induction1} this shows the added \texttt{pat} layer
rule is neither too strong nor too weak: it round-trips.
\end{lemma}
\begin{proof}
\uses{rule:isdefeq-pat}
A rule-by-rule induction in both cases. The \texttt{iota} branch adds the one-liner
\texttt{| pat hp hm he hr \_ ihall => exact .pat hp hm (hty he) hr ihall} at line 158, sending the
layer's rule to the core \texttt{VEnv.IsDefEq.pat} (\ref{rule:isdefeq-pat}) with the typing premise
pushed through \texttt{hty} and the side conditions supplied by the induction hypothesis.
Both declarations are in fact complete: their axiom sets are \texttt{propext} and
\texttt{Quot.sound} only, with no \texttt{sorryAx}. The \texttt{leanok} marker is withheld here
solely because the automated census attributes to this declaration range the commented-out
continuation of the file (lines 160--340), which does contain \texttt{sorry}s.
\end{proof}

\begin{definition}[One typing layer over an erased equality, \texttt{HasTypeU1}]
\label{inductive:expb-hastypeu1}
\lean{Lean4Lean.VEnv.HasTypeU1}
\leanok
\uses{struct:venv}
\srcloc{Lean4Lean/Experimental/StratifiedUntyped.lean}{17}
\thesisref{unique.tex \S\texttt{sec:unique}.}
As \texttt{HasType1}, except that the conversion rule appeals to an abstract \emph{type-erased}
equality \texttt{IsDefEqU1} $\Gamma\,A\,B$.
\end{definition}

\begin{definition}[One erased equality layer, \texttt{IsDefEqU1}, with the $\iota$ rule]
\label{inductive:expb-isdefequ1}
\lean{Lean4Lean.VEnv.IsDefEqU1, Lean4Lean.VEnv.IsDefEqU1.pat}
\leanok
\uses{inductive:expb-hastypeu1, struct:venv, def:pattern-check-realizes, inductive:pattern-matches}
\contrib{mixed}\srcloc{Lean4Lean/Experimental/StratifiedUntyped.lean}{33}
\thesisref{unique.tex \S\texttt{sec:unique} with the $\iota$ rule of \S\texttt{sec:kappa}.}
\texttt{IsDefEqU1} is the type-erased analogue of \texttt{IsDefEq1}: equivalence, congruences,
$\beta$, $\eta$, proof irrelevance and definitional axioms, all without a type index (and hence
without a \texttt{defeqDF} rule, so the \texttt{defEq} parameter is unused). The \texttt{iota} branch
adds \texttt{pat} (lines 51--55) with the same premises as the typed version, the side conditions
now being untyped equalities $\Gamma\vdash t.1\equiv t.2.1$.
\reviewnote{In the erased setting \texttt{Realizes} constrains only $t.1$ and $t.2.1$, so the third
component of every triple is unconstrained and the rule is a roundabout way of asserting
\texttt{r.2.OK}. Uniformity with the typed version is a defensible reason to keep the shape, but no
comment says so. Nothing in the file consumes the constructor: the would-be consumer
\texttt{IsDefEqU1.induction} is inside the commented-out block at lines 121--142.}
\end{definition}

\begin{lemma}[Decomposition into one erased layer]
\label{thm:expb-inductionu1}
\lean{Lean4Lean.VEnv.IsDefEq.inductionU1, Lean4Lean.VEnv.HasTypeU1.induction}
\leanok
\uses{inductive:expb-isdefequ1, inductive:expb-hastypeu1, struct:orderedstrong}
\contrib{mixed}\srcloc{Lean4Lean/Experimental/StratifiedUntyped.lean}{60}
\thesisref{unique.tex \S\texttt{sec:unique}.}
Under \texttt{OrderedStrong env} and a well-formed context, every $\Gamma\vdash e_1\equiv e_2 : A$
yields \texttt{HasTypeU1} derivations of both sides at $A$ and an \texttt{IsDefEqU1} derivation
relating them. Conversely \texttt{HasTypeU1.induction} (line 109, which keeps the weaker
\texttt{Ordered} hypothesis) collapses a layer back to \texttt{env.HasType}.
\end{lemma}
\begin{proof}
\uses{thm:isdefeq-strong}
\textbf{Not proved in Lean.} One \texttt{sorry} remains at line 77, the same \texttt{constDF} gap as
in \ref{thm:expb-induction1} and equally pre-existing. The rest goes through
\texttt{IsDefEq.strong} and an induction on the strong derivation; the \texttt{iota} branch adds the
\texttt{pat} case at lines 103--104, identical in shape to the typed version, and the same silent
\texttt{Ordered} to \texttt{OrderedStrong} strengthening at line 59.
\axfoot{direct sorry: StratifiedUntyped.lean:77 (constDF)}
\end{proof}

\section{Level-annotated judgments}

\begin{definition}[Level-annotated environment, \texttt{VEnv'}]
\label{structure:expb-venvprime}
\lean{Lean4Lean.VEnv', Lean4Lean.VEnv'.VConstant, Lean4Lean.VEnv'.VDefEq, Lean4Lean.VEnv'.empty, Lean4Lean.VEnv'.LE}
\leanok
\uses{struct:venv}
\srcloc{Lean4Lean/Experimental/Stronger.lean}{13}
\texttt{VEnv'.VConstant} and \texttt{VEnv'.VDefEq} extend the model's constant and definitional-axiom
records with the universe level of the type; \texttt{VEnv'} pairs a partial map of the former with a
predicate on the latter, and comes with an empty environment and an extension order
\texttt{VEnv'.LE} (reflexive and transitive).
\end{definition}

\begin{definition}[Erasure of a level-annotated environment, \texttt{VEnv'.out}]
\label{def:expb-venvprime-out}
\lean{Lean4Lean.VEnv'.out, Lean4Lean.VEnv'.empty_out}
\leanok
\uses{structure:expb-venvprime, struct:venv}
\contrib{mixed}\srcloc{Lean4Lean/Experimental/Stronger.lean}{25}
\texttt{VEnv'.out} forgets the level annotations: the constants are the underlying records, and the
definitional axioms are the images of the annotated ones. \texttt{empty\_out} says the empty
environment erases to the empty environment. The \texttt{iota} branch supplies the new field as
\texttt{pats \_ \_ := False} (line 28), so an erased environment registers no reduction rules.
\reviewnote{Sound but narrowing: every theorem in this file about \texttt{env.out} now speaks only
about $\iota$-free environments, so the experiment can never be reconnected to the $\iota$
development without giving \texttt{VEnv'} a real \texttt{pats} field. Given that the file dead-ends
in four \texttt{sorry}s the shortcut is pragmatic, but it changes the scope of the file's theorems
without a word of comment.}
\end{definition}

\begin{definition}[Definitional equality with explicit type levels, \texttt{IsDefEqStrong}]
\label{inductive:expb-isdefeqstrong-lvl}
\lean{Lean4Lean.VEnv'.IsDefEqStrong, Lean4Lean.VEnv'.Lookup, Lean4Lean.VEnv'.VCtx}
\leanok
\uses{structure:expb-venvprime}
\srcloc{Lean4Lean/Experimental/Stronger.lean}{58}
\thesisref{unique.tex \S\texttt{sec:unique}, aiming at the level clause of definitional inversion.}
Over a context \texttt{VCtx} of (type, level) pairs, \texttt{IsDefEqStrong} $\Gamma\,e_1\,e_2\,A\,u$
is a combined typing and equality judgment in which every judgment records the universe level $u$ of
its ascribed type $A$, every rule carries the sort-typing witnesses of the types it mentions, and an
extra \texttt{defeqL} rule converts along level equivalence.
\end{definition}

\begin{lemma}[Metatheory of the level-annotated judgment]
\label{family:expb-stronger-infra}
\lean{Lean4Lean.VEnv'.IsDefEqStrong.out, Lean4Lean.VEnv'.IsDefEqStrong.weakN, Lean4Lean.VEnv'.IsDefEqStrong.mono, Lean4Lean.VEnv'.IsDefEqStrong.weak0, Lean4Lean.VEnv'.IsDefEqStrong.instL, Lean4Lean.VEnv'.IsDefEqStrong.instN, Lean4Lean.VEnv'.IsDefEqStrong.defeqDF_l, Lean4Lean.VEnv'.IsDefEqStrong.isType', Lean4Lean.VEnv'.IsDefEqStrong.sort_invL, Lean4Lean.VEnv'.Ordered, Lean4Lean.VEnv'.CtxStrong, Lean4Lean.VEnv'.EnvStrong}
\leanok
\uses{inductive:expb-isdefeqstrong-lvl, def:expb-venvprime-out}
\srcloc{Lean4Lean/Experimental/Stronger.lean}{167}
The annotated judgment erases to the ordinary one (\texttt{out}), and is stable under weakening,
environment extension, universe instantiation and substitution; \texttt{isType'} is regularity and
\texttt{sort\_invL} pins the level index at a sort. \texttt{Ordered}, \texttt{CtxStrong} and
\texttt{EnvStrong} are the corresponding environment and context invariants.
\end{lemma}
\begin{proof}
\uses{def:expb-venvprime-out}
\leanok
Routine inductions mirroring the unannotated development in \texttt{Theory/Typing}, carrying the
level index along.
\end{proof}

\begin{lemma}[Product inversion in the level-annotated system]
\label{thm:expb-forallE-inv}
\lean{Lean4Lean.VEnv'.IsDefEqStrong.forallE_inv'}
\leanok
\uses{inductive:expb-isdefeqstrong-lvl, family:expb-stronger-infra}
\srcloc{Lean4Lean/Experimental/Stronger.lean}{478}
\thesisref{unique.tex, clause 2 of definitional inversion.}
Over an ordered environment satisfying \texttt{OnTypes (EnvStrong env)} and a strong context, if
either side of a judgment is $\forall A.\,B$, then $A$ is a type at some level $u$ and $B$ is a type
at some level $v$ in the context extended by $(A,u)$.
\end{lemma}
\begin{proof}
\uses{family:expb-stronger-infra}
\leanok
Induction on the derivation, generalised over the disjunction ``left side or right side is a
product''. The \texttt{extra} case is reduced to the environment's \texttt{EnvStrong} invariant, and
the $\beta$ case to the substitution lemma.
\end{proof}

\begin{lemma}[Level uniqueness (abandoned)]
\label{thm:expb-uniqL}
\lean{Lean4Lean.VEnv'.IsDefEqStrong.uniqL'}
\leanok
\uses{inductive:expb-isdefeqstrong-lvl, thm:expb-forallE-inv}
\srcloc{Lean4Lean/Experimental/Stronger.lean}{575}
\thesisref{unique.tex, the unique-typing theorem: this is an attempt to break its circularity by
proving the level component first.}
For two annotated judgments in a common context, if any of the four pairings of their subject terms
agree syntactically then their level indices are equivalent.
\end{lemma}
\begin{proof}
\textbf{Not proved in Lean.} Four \texttt{sorry}s remain at lines 624--628, one of them annotated
\texttt{looks like it needs unique typing :(}. The proof is a nested induction on both derivations;
it goes through for variables and constants but stalls in the \texttt{appDF}-versus-\texttt{appDF}
case, which is exactly the circularity the thesis warns about. Note also that the environment and
context hypotheses of the statement were commented out at line 574, so the theorem as stated is
weaker than the intended one.
\axfoot{direct sorry: Stronger.lean:624--628}
\end{proof}

\section{The SExpr calculus and unique typing}

\begin{definition}[Parameters of the \texttt{SExpr} development]
\label{class:expb-sexpr-params}
\lean{Lean4Lean.Params, Lean4Lean.Classification, Lean4Lean.Pattern.WF, Lean4Lean.SExpr.Params.extra_pat}
\leanok
\uses{struct:venv, inductive:pattern-core}
\srcloc{Lean4Lean/Experimental/SExpr.lean}{23}
\texttt{Params} is a class fixing an ordered environment, a universe count, an abstract family
\texttt{Pat p r} of reduction rules, and a \texttt{classify} map sending constant names to a
\texttt{Classification} (constructor, eta-constructor, symbol or inductive type, each with an
arity). Its hypotheses are that every rule's pattern is simple, is well formed for the
classification (\texttt{Pattern.WF}), and is unique among overlapping patterns. This is the
\texttt{SExpr}-side analogue of \texttt{VEnv.Params} (\ref{class:cr-params}).
\reviewnote{\texttt{Params.extra\_pat} is declared as a global \texttt{axiom} at line 614 rather than
as a class field, although the commented-out block at lines 36--44 shows it was meant to be one,
along with five further intended fields. Every result in this cluster is therefore conditional on a
global axiom that no instance is ever required to justify.}
\end{definition}

\begin{definition}[\texttt{SExpr} syntax and its substitution calculus]
\label{family:expb-sexpr-syntax}
\lean{Lean4Lean.SLevel, Lean4Lean.SExpr, Lean4Lean.SExpr.lift', Lean4Lean.SExpr.ClosedN, Lean4Lean.SExpr.Subst, Lean4Lean.SExpr.subst, Lean4Lean.SExpr.Subst.comp, Lean4Lean.SExpr.inst, Lean4Lean.SExpr.Skips, Lean4Lean.SExpr.Ctx.Lift', Lean4Lean.SExpr.Ctx.Inter}
\leanok
\uses{class:expb-sexpr-params}
\srcloc{Lean4Lean/Experimental/SExpr.lean}{49}
\thesisref{typesys.tex, the weakening and substitution lemmas.}
\texttt{SLevel} is the quotient of \texttt{VLevel} by semantic equality, chosen so that level
congruences disappear from the judgments. \texttt{SExpr} is a de Bruijn syntax over it with a
first-class type \texttt{Lift} of renamings, a substitution monoid \texttt{Subst} with composition,
identity, cons, tail and truncation, and the full set of interaction laws between lifting,
substitution and level instantiation. \texttt{Skips} records which variables a term avoids, and
\texttt{Ctx.Lift'} and \texttt{Ctx.Inter} describe liftings of contexts and how two liftings of a
common context intersect. This API is cross-referenced from the shipping code, at
\srcloc{Lean4Lean/Theory/VExpr.lean}{836}, so the file is not entirely dead.
\end{definition}

\begin{definition}[The \texttt{SExpr} weak definitional equality]
\label{inductive:expb-sexpr-isdefeq}
\lean{Lean4Lean.SExpr.IsDefEq, Lean4Lean.SExpr.IsDefEqStrong, Lean4Lean.SExpr.IsDefEq.strong, Lean4Lean.SExpr.IsDefEqStrong.defeq}
\leanok
\uses{class:expb-sexpr-params, family:expb-sexpr-syntax}
\srcloc{Lean4Lean/Experimental/SExpr.lean}{590}
\thesisref{typesys.tex \S2, the definitional equality judgment; unique.tex \S\texttt{sec:unique}.}
$\Gamma\vdash e_1\equiv e_2 : A$ is the usual typed equality over \texttt{SExpr} with one extra rule,
a heterogeneous transitivity \texttt{trans'} letting the middle term of a chain of type equalities be
typed at a different sort. \texttt{IsDefEqStrong} is the variant carrying subject-reduction witnesses
inline, including a \texttt{const} rule that demands a \texttt{CtorBundle} exhibiting each constant's
type as a telescope ending in an inductive head. The two bridges \texttt{IsDefEq.strong} and
\texttt{IsDefEqStrong.defeq} between them are stated at lines 679--680 and are literally
\texttt{:= sorry}.
\reviewnote{This judgment has \emph{no} $\iota$ or \texttt{pat} rule: the pattern-based rule is
commented out at lines 609--610, and reduction rules enter only through the \texttt{extra} rule over
\texttt{env.defeqs} plus the global \texttt{Params.extra\_pat} axiom. The \texttt{iota} branch
extended \texttt{VEnv.IsDefEq} and \texttt{VEnv.Params} but left this parallel judgment untouched,
so the two halves of the project now model $\iota$ differently.}
\axfoot{direct sorry: SExpr.lean:679 (IsDefEq.strong), :680 (IsDefEqStrong.defeq)}
\end{definition}

\begin{definition}[Operational judgments over \texttt{SExpr}]
\label{family:expb-sexpr-operational}
\lean{Lean4Lean.SExpr.WHRed, Lean4Lean.SExpr.WHRedS, Lean4Lean.SExpr.WHNF, Lean4Lean.SExpr.ParRed, Lean4Lean.SExpr.ParRedS, Lean4Lean.SExpr.InferType, Lean4Lean.SExpr.InferTypeS, Lean4Lean.SExpr.NormalEq, Lean4Lean.SExpr.CRDefEq, Lean4Lean.SExpr.Ctx.Subst, Lean4Lean.SExpr.Ctx.SubstEq}
\leanok
\uses{inductive:expb-sexpr-isdefeq, family:expb-sexpr-syntax}
\srcloc{Lean4Lean/Experimental/SExpr.lean}{947}
\thesisref{typesys.tex, the reduction and algorithmic-typing items; unique.tex
\S\texttt{sec:church\_rosser}.}
Weak-head reduction \texttt{WHRed} (head application, $\beta$, and an \texttt{extra} rule that fires
an abstract \texttt{Pat} rule whose side conditions are discharged by a list of definitional
equalities) closes to \texttt{WHRedS}; \texttt{WHNF} is its normal-form predicate.
\texttt{WHRed.determ} establishes determinism only for the $\beta$/application fragment: all five
case splits where an \texttt{extra} step meets anything, including another \texttt{extra} step, are
\texttt{sorry} (lines 987, 992, 995--997). \texttt{ParRed} is the parallel version with the same \texttt{extra}
rule. \texttt{InferType} is a syntax-directed inference relation, deterministic and stable under
weakening and substitution, with an up-to-reduction variant \texttt{InferTypeS}. \texttt{NormalEq}
and \texttt{CRDefEq} repeat the proof-irrelevance and Church-Rosser relations of
Section~\ref{chap:experimental-reduction} in typed, \texttt{SExpr} form.
\reviewnote{The soundness bridges \texttt{WHRedS.defeq} (line 1008), \texttt{InferType.hasType}
(1091), \texttt{InferTypeS.hasType} (1167), \texttt{WHRedS.parRedS} (1182) and \texttt{Ctx.Subst.id}
(770) are all \texttt{sorry}, as is \texttt{CRDefEq.trans} (1295), so the whole operational layer is
stated but not connected to the declarative one. Note that the pattern rule that
\texttt{SExpr.IsDefEq} lacks \emph{is} present here, in \texttt{WHRed} and \texttt{ParRed}: the two
layers disagree about how reduction rules enter.}
\end{definition}

\begin{definition}[Bundled \texttt{SExpr} typing judgment, \texttt{HasTypeS}]
\label{inductive:expb-hastypes}
\lean{Lean4Lean.SExpr.HasTypeS}
\leanok
\uses{inductive:expb-sexpr-isdefeq}
\srcloc{Lean4Lean/Experimental/UniqueTyping.lean}{29}
\thesisref{unique.tex \S\texttt{sec:unique}, the judgment used in the proof of unique typing ``which
has no weakening rule''.}
\texttt{HasTypeS} $\Gamma\,e\,A\,b$ is a typing judgment indexed by a boolean: at $b = $
\texttt{false} the derivation is structural, exactly one syntax-directed rule with the sort
witnesses of all types carried inline; at $b = $ \texttt{true} the conversion rule is additionally
available, and \texttt{base} injects the structural fragment into the conversion-closed one. Its
docstring correctly calls it the \texttt{SExpr} analogue of \texttt{HasTypeStrong} on the
\texttt{VExpr} side, stripped of the stratification index.
\end{definition}

\begin{theorem}[Type uniqueness for the bundled judgment]
\label{thm:expb-hastypes-uniq}
\lean{Lean4Lean.SExpr.HasTypeS.uniq, Lean4Lean.SExpr.HasTypeS.toStructural, Lean4Lean.SExpr.HasTypeS.unfold, Lean4Lean.SExpr.HasTypeS.hasType}
\leanok
\uses{inductive:expb-hastypes, inductive:expb-sexpr-isdefeq}
\srcloc{Lean4Lean/Experimental/UniqueTyping.lean}{92}
\thesisref{unique.tex, the unique-typing theorem.}
Every bundled derivation projects to an \texttt{IsDefEq} derivation (\texttt{hasType}) and factors
as a structural derivation followed by a chain of conversions (\texttt{toStructural},
\texttt{unfold}). Consequently, for any $b_1, b_2$, from \texttt{HasTypeS} $\Gamma\,e\,A\,b_1$ and
\texttt{HasTypeS} $\Gamma\,e\,B\,b_2$ there is $u$ with $\Gamma\vdash A\equiv B : \mathtt{.sort}\ u$.
\end{theorem}
\begin{proof}
\uses{thm:expa-forallE-inv, thm:expa-sort-inv}
\leanok
Induction on the first derivation, reducing the second to structural form at each step and matching
the two syntax-directed rules. The application case inverts the two product types with
\texttt{SExpr.forallE\_inv} (\ref{thm:expa-forallE-inv}) and substitutes with
\texttt{Ctx.SubstEq}; the sort and product cases use \texttt{SExpr.sort\_inv}
(\ref{thm:expa-sort-inv}); the conversion case appeals to the heterogeneous \texttt{trans'}. This is
the closest thing in the whole experimental tree to the thesis's unique-typing theorem, and it is
genuinely proved from its hypotheses.
\axfoot{sorryAx via the ShapeLogRelAdequacy inversions; Lean4Lean.SExpr.Params.extra\_pat}
\end{proof}

\begin{theorem}[Sort uniqueness for the \texttt{SExpr} equality]
\label{thm:expb-uniq-sort}
\lean{Lean4Lean.SExpr.IsDefEq.toHasTypeS, Lean4Lean.SExpr.IsDefEq.uniq_sort}
\leanok
\uses{thm:expb-hastypes-uniq, inductive:expb-hastypes, inductive:expb-sexpr-isdefeq}
\srcloc{Lean4Lean/Experimental/UniqueTyping.lean}{138}
\thesisref{unique.tex, clause 1 of definitional inversion.}
\texttt{toHasTypeS} turns any $\Gamma\vdash e_1\equiv e_2 : A$ into bundled typings of both sides at
$A$. Hence if $\Gamma\vdash e_1\equiv e_2 : \mathtt{.sort}\ u$ and
$\Gamma\vdash e_2\equiv e_3 : \mathtt{.sort}\ v$ then $u = v$.
\end{theorem}
\begin{proof}
\uses{thm:expb-hastypes-uniq}
\leanok
\texttt{toHasTypeS} opens with \texttt{replace h := h.strong} and then replays the strong rules; the
\texttt{trans'} case is where the extra power is needed, and it is closed by \texttt{HasTypeS.uniq}
plus \texttt{sort\_inv}. \texttt{uniq\_sort} then applies \texttt{HasTypeS.uniq} to the two typings
of the middle term and inverts.
\reviewnote{\texttt{UniqueTyping.lean} contains no textual \texttt{sorry}, yet
\texttt{IsDefEq.strong} is \texttt{:= sorry} at \srcloc{Lean4Lean/Experimental/SExpr.lean}{679}, so
this theorem and everything after it in the file are conditional on an admitted statement. The
docstring on \texttt{toHasTypeS} is candid about this, but any status derived from grepping for
\texttt{sorry} would mislabel the file as complete.}
\axfoot{sorryAx via SExpr.IsDefEq.strong (SExpr.lean:679); Lean4Lean.SExpr.Params.extra\_pat}
\end{proof}

\begin{lemma}[Heterogeneous transitivity is admissible]
\label{thm:expb-isdefeq-prime}
\lean{Lean4Lean.SExpr.IsDefEq', Lean4Lean.SExpr.IsDefEq'.toIsDefEq, Lean4Lean.SExpr.IsDefEq.toIsDefEq', Lean4Lean.SExpr.IsDefEq.iff_isDefEq'}
\leanok
\uses{inductive:expb-sexpr-isdefeq, thm:expb-uniq-sort}
\srcloc{Lean4Lean/Experimental/UniqueTyping.lean}{195}
\thesisref{unique.tex \S\texttt{sec:unique}.}
\texttt{IsDefEq'} is \texttt{IsDefEq} with the \texttt{trans'} rule removed. The two systems derive
exactly the same judgments: \texttt{iff\_isDefEq'}.
\end{lemma}
\begin{proof}
\uses{thm:expb-uniq-sort}
\leanok
The forward direction is a rule-for-rule embedding. The backward direction is an induction in which
the \texttt{trans'} case is discharged by \texttt{uniq\_sort}, which shows the two sorts coincide, so
that ordinary \texttt{trans} applies.
\reviewnote{Neither \texttt{IsDefEq} nor \texttt{IsDefEq'} has an $\iota$ rule, so this equivalence
has not been checked against the system the \texttt{iota} branch actually extended.}
\axfoot{sorryAx via SExpr.IsDefEq.strong (SExpr.lean:679); Lean4Lean.SExpr.Params.extra\_pat}
\end{proof}

\section{Unfinished semantic models}

\begin{definition}[Axiomatized semantic domain (Coquand-Huber)]
\label{family:expb-thierry-domain}
\lean{D, DF, FinMut, FinHasType, El, Expr, Expr.eval, WHRedTS}
\leanok
\srcloc{Lean4Lean/Experimental/Thierry.lean}{11}
\thesisref{no counterpart: normalization.tex is a stub ending in \texttt{UNFINISHED}.}
A partial formalization of Coquand and Huber, \emph{An Adequacy Theorem for Dependent Type Theory},
written (per its own header) to check the author's understanding of the monotonicity theorem. The
semantic domain \texttt{D} is an \texttt{axiom}, as are its order, bottom, universe, $\lambda$ and
$\Pi$ constructors and application; \texttt{DF} is the subtype of well-behaved endofunctions.
\texttt{FinMut} is a datatype of finite approximations with an axiomatized embedding into
\texttt{D}, \texttt{FinHasType} is a shape-level typing relation, \texttt{El} reads a semantic type
back as a predicate, and \texttt{Expr} with \texttt{Expr.eval} and the typed reduction
\texttt{WHRedTS} is the object syntax. These declarations live in the root namespace, so several of
them collide by name with the corresponding ones in \texttt{Thierry2.lean}.
\reviewnote{Line 9 declares \texttt{axiom mySorry : $\alpha$} with an auto-bound implicit, that is,
an inhabitant of every type and in particular a proof of \texttt{False}. It is used in
\texttt{DF.comp}, \texttt{DF.bot} and elsewhere, so this file is deliberately inconsistent. Nothing
imports it, but nothing downstream may ever be allowed to.}
\axfoot{the axiom mySorry (Thierry.lean:9) proves False; 14 further sorry}
\end{definition}

\begin{lemma}[Monotonicity of the adequacy interpretation]
\label{family:expb-thierry-mono}
\lean{HasTypeF, HasTypeFamF, HasPiFamF, HasTypeF.mono, HasPiFamF.mono, HasTypeFamF.mono, HasTypeF.mono_r, HasPiFamF.mono_r}
\leanok
\uses{family:expb-thierry-domain}
\srcloc{Lean4Lean/Experimental/Thierry.lean}{194}
The mutually recursive relations \texttt{HasTypeF} (a term is adequately interpreted at a shape),
\texttt{HasTypeFamF} (a type family) and \texttt{HasPiFamF} (a $\Pi$ family) are monotone in their
shape arguments: covariantly in the term shape (\texttt{mono}, at line 248) and contravariantly in
the type shape (\texttt{mono\_r}, at line 321). This is the one substantive result the file set out
to check.
\end{lemma}
\begin{proof}
\uses{family:expb-thierry-domain}
\leanok
A large mutual induction on the shape measure, terminating by \texttt{meas\_apply\_lt}, with case
analysis on the shape of the type. The \texttt{mono} theorems themselves are complete, but they rest
on a dozen admitted auxiliaries (\texttt{FinMut.LE.rfl}, \texttt{FinHasType.mono},
\texttt{FinHasType.lem5}, \texttt{El\_iff}, \texttt{meas\_apply\_lt}, \texttt{subst} and others) and
on the inconsistent \texttt{mySorry} axiom, so they certify understanding, not truth.
\axfoot{the axiom mySorry (Thierry.lean:9); sorryAx via FinHasType.mono, El\_iff, meas\_apply\_lt, subst}
\end{proof}

\begin{definition}[The finite shape approximation lattice]
\label{family:expb-thierry2-shape}
\lean{Shape, ShapeFun, Shape.LE, Shape.Join, Shape.Compat, Shape.app, Shape.lift, Shape.HasType, Shape.HasTypeU}
\leanok
\srcloc{Lean4Lean/Experimental/Thierry2.lean}{21}
\thesisref{no counterpart (normalization.tex is a stub).}
A second, more self-contained pass at the same adequacy argument, in which the abstract domain is
replaced by a concrete construction. \texttt{Shape n} is an $n$-level approximation of a semantic
value (bottom, the universe, a $\Pi$ given by a domain shape and a finite table of
argument/result shapes, or a $\lambda$ given by such a table), with a decidable order, a
compatibility relation, binary joins, application, lifting between levels, and a shape-level typing
relation characterised inductively by \texttt{HasTypeU}.
\reviewnote{The order-theoretic core is assumed, not proved: \texttt{Shape.Compat.def},
\texttt{Shape.Join.mk}, \texttt{ShapeFun.Join.mk}, \texttt{Shape.HasType.mono},
\texttt{ShapeFun.LE.bot}, \texttt{ShapeFun.bot\_app} and \texttt{Shape.app\_mono\_r} are all
\texttt{sorry}. The same construction reappears in \texttt{MoreStepIndexed.lean} and in
\texttt{ShapeLogRel.lean}: three near-copies with three independently admitted sets of lattice laws.}
\end{definition}

\begin{lemma}[Concrete domain from shapes, and monotonicity of the interpretation]
\label{family:expb-thierry2-model}
\lean{D, DF, D.lam, D.unlam, D.pi, Shape.embed, interp, DefEqF, DefEqF.mono, DefEqF.mono_r, DefEqPiF.mono, DefEqLamF.mono}
\leanok
\uses{family:expb-thierry2-shape}
\srcloc{Lean4Lean/Experimental/Thierry2.lean}{351}
A semantic value \texttt{D} is a downward-closed, join-complete set of shapes; \texttt{DF} is the
corresponding function space, with $\lambda$, its inverse \texttt{unlam} and $\Pi$ constructed on
top. \texttt{interp} evaluates the object syntax into it, and \texttt{DefEqF} interprets a
definitional-equality judgment at a shape. The payload is the monotonicity of \texttt{DefEqF} and of
its $\Pi$ and $\lambda$ variants in both shape arguments.
\end{lemma}
\begin{proof}
\uses{family:expb-thierry2-shape}
\textbf{Not proved in Lean.} Two of the declarations this family bundles are literal
\texttt{sorry}s, \texttt{D.pi} (line 472) and \texttt{DF.mk'} (line 421); neither name occurs in the
\texttt{DefEqF}/\texttt{DefEqPiF}/\texttt{DefEqLamF} \texttt{mono}/\texttt{mono\_r} theorems' own
bodies, which instead go through mutual induction on shape levels, using the join and compatibility
lemmas of the shape lattice. Those lattice lemmas are themselves partly \texttt{sorry}, so the family
is transitively tainted regardless of the textual split.
\reviewnote{\texttt{DF.mk'} (line 421) and \texttt{D.pi} (line 472) are \texttt{sorry}, so the model
is incomplete at exactly the connective the adequacy theorem is about; the judgment being
interpreted, \texttt{DefEq}, is itself an \texttt{axiom} (lines 607--609) rather than a definition;
and this file too opens with \texttt{axiom mySorry}.}
\axfoot{the axiom mySorry (Thierry2.lean:9); sorryAx via D.pi, DF.mk', Shape.Join.mk, Shape.Compat.def}
\end{proof}

\begin{definition}[Semantic type classifiers over \texttt{SExpr} (sketch)]
\label{family:expb-stepindexed}
\lean{Lean4Lean.SExpr.Classifier', Lean4Lean.SExpr.Classifier.HasTy, Lean4Lean.SExpr.NormalType, Lean4Lean.SExpr.normalize, Lean4Lean.SExpr.normalize_prop, Lean4Lean.SExpr.IsTy', Lean4Lean.SExpr.IsTy, Lean4Lean.SExpr.IsTyN, Lean4Lean.SExpr.IsTy.def}
\leanok
\uses{family:expb-sexpr-operational}
\srcloc{Lean4Lean/Experimental/StepIndexed.lean}{8}
\thesisref{normalization.tex, the unwritten strong-normalization section.}
A \texttt{Classifier} pairs a level with a membership predicate on terms; \texttt{normalize}
$\Gamma\,A$ classically picks a weak-head normal \emph{type} reachable from $A$, if one exists, and
\texttt{normalize\_prop} extracts its defining property. \texttt{IsTy'} is a coinductive rendering
of ``$A$ is a type classified by $C$'' with three cases (sort, product, stuck).
\reviewnote{\texttt{IsTy}, \texttt{IsTyN} and the unfolding equation \texttt{IsTy.def} are
\texttt{axiom}s: the fixpoint the file sets out to construct is posited rather than built, which
begs the file's own question. \texttt{IsTyN} is declared and never used, and roughly 40\% of the
88-line file is commented out. This is scratch, not a result.}
\end{definition}

\begin{definition}[Step-indexed propositions and bounded judgments]
\label{family:expb-morestepindexed}
\lean{Lean4Lean.SExpr.IProp, Lean4Lean.SExpr.BIProp, Lean4Lean.SExpr.WHRedUpToN, Lean4Lean.SExpr.ParRedN, Lean4Lean.SExpr.InferTypeN, Lean4Lean.SExpr.NormalEqN, Lean4Lean.SExpr.CRDefEqN, Lean4Lean.SExpr.CheckTypeN}
\leanok
\srcloc{Lean4Lean/Experimental/MoreStepIndexed.lean}{15}
\thesisref{normalization.tex (unwritten).}
\texttt{IProp} is an upward-closed family of propositions indexed by a step count and \texttt{BIProp
n} its bounded version, with restriction operators and the expected idempotence. On top of them the
$n$-bounded reduction, parallel-reduction, inference and normal-equality judgments are combined into
\texttt{CRDefEqN} and \texttt{CheckTypeN}, each with a monotonicity lemma.
\reviewnote{The four bounded judgments and their four monotonicity lemmas are \texttt{axiom}s
(lines 38--45): the entire operational layer this file is about is assumed. The step-index
bookkeeping is the only content that is actually proved.}
\end{definition}

\begin{definition}[Shape-indexed logical relation]
\label{family:expb-morestepindexed-logrel}
\lean{Lean4Lean.SExpr.Classifier', Lean4Lean.SExpr.Classifier.DefEq, Lean4Lean.SExpr.LogRel, Lean4Lean.SExpr.LogRel.TypeEq, Lean4Lean.SExpr.LogRelK.bot, Lean4Lean.SExpr.LogRelK.sort, Lean4Lean.SExpr.LogRelK.forallE}
\leanok
\uses{family:expb-morestepindexed, family:expb-thierry2-shape}
\srcloc{Lean4Lean/Experimental/MoreStepIndexed.lean}{338}
\thesisref{normalization.tex (unwritten).}
A \texttt{Classifier} assigns to a pair of types a family of equality predicates on terms, monotone
in the shape; \texttt{LogRel} $\Gamma\,A\,B\,n$ assigns such a classifier to each shape at level
$n$, with \texttt{LogRel.TypeEq} the membership predicate; \texttt{LogRelK.bot},
\texttt{LogRelK.sort} and \texttt{LogRelK.forallE} build the classifiers for the two type formers
out of the previous level.
\reviewnote{The file contains \texttt{\#exit} at line 420, so the two declarations that would
assemble one level from the previous and iterate it, \texttt{TypeEqS} (421) and \texttt{TypeEq}
(443), are never elaborated and do not exist in the compiled environment; neither does
\texttt{TypeEq.mono}, which is commented out at line 489. Of the five textual \texttt{sorry}s in the
file only one, \texttt{ShapeFun.app\_mono\_l} at line 310, is live; the rest are inside comments.
Lines 63--336 duplicate the \texttt{Shape} lattice of \texttt{Thierry2.lean} wholesale, admitted
order laws included.}
\axfoot{sorryAx via ShapeFun.app\_mono\_l (MoreStepIndexed.lean:310)}
\end{definition}

\begin{definition}[Inverse-limit domain for the \texttt{SExpr} syntax functor]
\label{family:expb-domaintheory}
\lean{Lean4Lean.SExpr.SExprF, Lean4Lean.SExpr.SExprF.map, Lean4Lean.SExpr.SExprF.select, Lean4Lean.SExpr.DomN, Lean4Lean.SExpr.DomN.up, Lean4Lean.SExpr.DomN.down, Lean4Lean.SExpr.DomN.cast, Lean4Lean.SExpr.Dom, Lean4Lean.SExpr.Dom.in, Lean4Lean.SExpr.Dom.out}
\leanok
\uses{family:expb-sexpr-syntax}
\srcloc{Lean4Lean/Experimental/DomainTheory.lean}{8}
\thesisref{normalization.tex (unwritten).}
\texttt{SExprF V F} is one layer of the \texttt{SExpr} syntax with abstract slots for subterms and
for binders, with a functorial \texttt{map} and a dependent eliminator \texttt{select};
\texttt{DomN n} is its $n$-fold unfolding, equipped with transfer maps \texttt{up}, \texttt{down} and
\texttt{cast} whose coherence laws are proved by induction on the index; \texttt{Dom} is the inverse
limit, the coherent sequences of approximations. \texttt{Dom.in} folds one layer into the limit.
\reviewnote{\texttt{Dom.out} (line 211), which should unfold it again, contains the \texttt{stop}
tactic at line 223, and \texttt{stop} is shorthand for \texttt{repeat sorry}: the definition is
\texttt{sorryAx}-backed even though a text search for \texttt{sorry} finds nothing in this file. The
file then ends at \texttt{Dom.fin} (line 240), a stub mapping every finite element to bottom.}
\axfoot{sorryAx via the stop tactic in Dom.out (DomainTheory.lean:223)}
\end{definition}

\section{Contribution summary}

The \texttt{iota} branch changed 28 lines in this group and \texttt{trproj} changed none. In
dependency order the contributions are: the interface field \texttt{Typing.pat\_env}
(\ref{structure:expb-typing}, one line plus a supporting import), which is what lets the two new
cases of
\ref{thm:expb-todyping} and \ref{thm:expb-isdefequ-church-rosser} (eight lines together) promote an
environment-registered rule to an abstract \texttt{Pat} rule; the layer constructor
\texttt{IsDefEq1.pat} with the two cases that produce and consume it
(\ref{inductive:expb-isdefeq1}, \ref{thm:expb-induction1}, \ref{thm:expb-isdefeq1-induction}, nine
lines) and its type-erased twin (\ref{inductive:expb-isdefequ1}, \ref{thm:expb-inductionu1}, eight
lines); and the stub \texttt{pats \_ \_ := False} in \ref{def:expb-venvprime-out} (one line). Two of
those eight-line groups also carry a silent strengthening of an exported hypothesis from
\texttt{Ordered} to \texttt{OrderedStrong}.

Why they exist: \texttt{VEnv} grew a \texttt{pats} field and \texttt{VEnv.IsDefEq} grew a
\texttt{pat} constructor, so every inductive over \texttt{IsDefEq} in the tree acquired a new case
and every construction of a \texttt{VEnv} acquired a new field. CI builds
\texttt{Lean4Lean.Experimental} even though \texttt{lake} does not, so the experiments had to be
repaired rather than left broken.

How much of it is verification work and how much is compile-keeping: only the nine lines in
\texttt{Stratified.lean} do something that nothing else in the repository does, namely decide how
the $\iota$ rule stratifies and check that the decision round-trips. The eight lines in
\texttt{NormalEq.lean} and \texttt{ParallelReduction.lean} are the same argument the branch already
wrote in \texttt{Theory/Typing/ChurchRosser.lean} (\ref{class:cr-params},
\ref{axiom:cr-params-pat-env}, \ref{thm:isdefeq-church-rosser}), in files whose headers ask for their
own deletion. The eight in \texttt{StratifiedUntyped.lean} are dead, since the only consumer is
commented out. The one in \texttt{Stronger.lean} narrows the file to $\iota$-free environments.

What remains open. In this group, all of it: the two theorems the branch extended,
\ref{thm:expb-isdefequ-church-rosser} and \ref{thm:expb-induction1}, are conditional on pre-existing
master \texttt{sorry}s, so nothing here is a proved result. The one substantive gap the branch left
is the \texttt{SExpr} side: \ref{inductive:expb-sexpr-isdefeq} still has no $\iota$ rule, so
\ref{thm:expb-hastypes-uniq}, \ref{thm:expb-uniq-sort} and \ref{thm:expb-isdefeq-prime} are results
about a system that does not model the reduction rules the rest of the project now has.

\section{Review notes}

\begin{itemize}
\item \textbf{Medium.} \srcloc{Lean4Lean/Experimental/ParallelReduction.lean}{909}: the $\iota$
case added here, together with the \texttt{Typing.pat\_env} field at
\srcloc{Lean4Lean/Experimental/NormalEq.lean}{97}, duplicates exactly what the branch added to
\texttt{Theory/Typing/ChurchRosser.lean} (the \texttt{pat} case of \texttt{IsDefEq.church\_rosser},
lines 1390--1394, and the \texttt{Params.pat\_env} field). Both files carry an upstream
\texttt{TODO: remove, this is now part of ChurchRosser.lean}. The duplication is forced by CI, but
the $\iota$ rule is now maintained in two places with no cross-reference comment in either. The
other $\iota$ case in this file, \texttt{VEnv.IsDefEq.toTyping} at
\srcloc{Lean4Lean/Experimental/ParallelReduction.lean}{857} (\ref{thm:expb-todyping}), is not a
duplicate of anything: it bridges a concrete derivation into the abstract \texttt{Typing} record
that \texttt{ChurchRosser.lean}, working directly over a \texttt{[Params]} instance, has no need for.
\item \textbf{Medium.} \srcloc{Lean4Lean/Experimental/SExpr.lean}{981}: the chapter text describes
\texttt{WHRed} as deterministic (\texttt{WHRed.determ}), but that theorem's proof is \texttt{sorry}
in all five case splits where an \texttt{extra} (pattern-rule) step meets anything, including another
\texttt{extra} step (lines 987, 992, 995--997); only the $\beta$/application fragment is actually
proved deterministic.
\item \textbf{Medium.} \srcloc{Lean4Lean/Experimental/Stratified.lean}{79} and
\srcloc{Lean4Lean/Experimental/StratifiedUntyped.lean}{59}: the hypothesis of an exported theorem
was silently strengthened from \texttt{Ordered env} to \texttt{OrderedStrong env}. The change is
forced and loses no realistic instance, but neither site explains it.
\item \textbf{Medium.} \srcloc{Lean4Lean/Experimental/SExpr.lean}{609}: the whole \texttt{SExpr}
line of work still has no $\iota$ rule in its declarative equality. The pattern-based rule is
commented out and reduction rules enter only through \texttt{env.defeqs} plus the global
\texttt{axiom Params.extra\_pat} at line 614, while \texttt{WHRed} and \texttt{ParRed} on the same
file \emph{do} carry a \texttt{Pat}-based rule. The branch extended \texttt{VEnv.IsDefEq} and
\texttt{VEnv.Params} and left this parallel development untouched.
\item \textbf{Medium.} \srcloc{Lean4Lean/Experimental/UniqueTyping.lean}{138}: the file contains no
textual \texttt{sorry} but \texttt{IsDefEq.toHasTypeS} opens with \texttt{h.strong}, and
\texttt{SExpr.IsDefEq.strong} is \texttt{:= sorry}. \texttt{uniq\_sort} (172),
\texttt{toIsDefEq'} (240) and \texttt{iff\_isDefEq'} (260) are all conditional; a status derived
from grepping would mislabel the file as complete.
\item \textbf{Medium.} \srcloc{Lean4Lean/Experimental/DomainTheory.lean}{223}: \texttt{Dom.out} is
\texttt{sorryAx}-backed through the \texttt{stop} tactic, which is \texttt{repeat sorry}, again
despite no literal \texttt{sorry} in the file. The same pitfall occurs at
\srcloc{Lean4Lean/Experimental/LogRel.lean}{320} and
\srcloc{Lean4Lean/Experimental/ShapeLogRel.lean}{1703}.
\item \textbf{Medium.} \srcloc{Lean4Lean/Experimental/MoreStepIndexed.lean}{420}: a
\texttt{\#exit} truncates elaboration, so \texttt{TypeEqS} (421) and \texttt{TypeEq} (443), the two
declarations the file exists to build, are never compiled and do not appear in the environment. A
reader of the source cannot tell this from the text, and the file still builds green.
\item \textbf{High.} \srcloc{Lean4Lean/Experimental/Thierry.lean}{9} and
\srcloc{Lean4Lean/Experimental/Thierry2.lean}{9}: \texttt{axiom mySorry : $\alpha$} is an inhabitant
of every type and therefore a proof of \texttt{False}. It is used in \texttt{DF.comp},
\texttt{DF.bot} and elsewhere. These files are scratch formalizations of an external paper and are
imported by nothing, but the axiom must be recorded so that nothing downstream is ever allowed to
depend on them. \texttt{Thierry2.lean} additionally axiomatizes the judgment it interprets
(\texttt{DefEq}, \texttt{DefEq.U}, \texttt{DefEq.left}, lines 607--609).
\item \textbf{Low.} \srcloc{Lean4Lean/Experimental/NormalEq.lean}{97}: \texttt{Typing.pat\_env} has
no docstring, unlike the field it mirrors in \texttt{ChurchRosser.lean}, which documents it. The
\texttt{Typing} structure is never instantiated anywhere in the repository, so the added obligation
is never discharged and gives no evidence that the $\iota$ rule is realisable in this interface.
\item \textbf{Low.} \srcloc{Lean4Lean/Experimental/StratifiedUntyped.lean}{51}: the added
\texttt{IsDefEqU1.pat} constructor is never eliminated, since the only consumer is inside the
commented-out block at lines 121--142, and its \texttt{chk} list has an entirely unconstrained third
component in the erased setting where \texttt{Check.OK} would say the same thing directly.
\item \textbf{Low.} \srcloc{Lean4Lean/Experimental/Stronger.lean}{28}: \texttt{pats \_ \_ := False}
silently restricts every theorem about erased environments in the file to $\iota$-free ones. Adding
a real \texttt{pats} field to \texttt{VEnv'} was the general fix; the file dead-ends in four
\texttt{sorry}s anyway, so the shortcut is pragmatic rather than wrong.
\item \textbf{Low.} \srcloc{Lean4Lean/Experimental/ParallelReduction.lean}{699}: the two
pre-existing master \texttt{sorry}s in \texttt{NormalEq.parRed} (699, 718), precisely the cases where
a pattern reduction meets a proof-irrelevance step, make \ref{thm:expb-parred-church-rosser},
\ref{family:expb-crdefeq} and \ref{thm:expb-isdefequ-church-rosser} conditional in this file. The
$\iota$ case the branch added therefore extends an unproved theorem; the corresponding case in
\texttt{ChurchRosser.lean} is the one that matters.
\item \textbf{Low.} \srcloc{Lean4Lean/Experimental/MoreStepIndexed.lean}{63}: the
\texttt{Shape}/\texttt{ShapeFun} approximation lattice exists in three near-identical copies, here,
at \srcloc{Lean4Lean/Experimental/Thierry2.lean}{21} and in \texttt{ShapeLogRel.lean}, each with its
own partly admitted order laws and nothing keeping them in sync.
\item \textbf{Low.} \srcloc{Lean4Lean/Experimental/StepIndexed.lean}{59}: \texttt{IsTy},
\texttt{IsTyN} and \texttt{IsTy.def} are \texttt{axiom}s positing the existence and the unfolding
equation of the very fixpoint the file sets out to construct; \texttt{IsTyN} is declared and never
used, and about 40\% of the file is commented out.
\item \textbf{Low.} \srcloc{Lean4Lean/Experimental/NormalEq.lean}{145}: several declarations here
carry fully qualified names that already exist in \texttt{Theory/Typing/ChurchRosser.lean}, for
instance \texttt{Lean4Lean.Pattern.Matches.hasType} and \texttt{Lean4Lean.Pattern.Check.OK.weakN}.
The two modules are never imported together so Lean accepts it, but any name-indexed tooling (an
axiom census, a blueprint cross-reference, a documentation generator) will conflate the sorry-free
copy here with the \texttt{sorryAx}-tainted one there.
\item \textbf{Low.} \srcloc{Lean4Lean/Experimental/Stratified.lean}{160} and
\srcloc{Lean4Lean/Experimental/StratifiedUntyped.lean}{121}: two thirds of each file is a
commented-out block of abandoned unique-typing and weakening-inversion attempts, containing further
\texttt{sorry}s. Besides being dead weight, it makes the last live declaration of each file look
admitted to any line-range-based analysis.
\end{itemize}
